% =============================================================
% Chapter 2: Deep Learning Pioneers and Godfathers
% Book: AI Figures — Who's Who in the Age of AI
% =============================================================

\chapter{Deep Learning Pioneers and Godfathers}
\label{ch:pioneers}

% Chapter overview: The individuals who built the intellectual and engineering
% foundations of modern deep learning — from backpropagation and CNNs to
% Transformers and reinforcement learning.  Their stories reveal how a
% marginalized research agenda became the dominant paradigm of 21st-century
% technology.

如果说人工智能是一座正在拔地而起的摩天大楼，那么本章讲述的就是那些打地基的人。
从1980年代在学术界边缘坚持神经网络研究的``异端''，到2017年在Google内部
悄然写出改变世界论文的八位工程师，再到用AI解决蛋白质折叠问题并获得诺贝尔奖的
科学家——这些人的故事构成了现代AI最激动人心的编年史。

他们中有些人等待了三十年才看到自己的理论被验证（Hinton、LeCun、Bengio），
有些人在二十多岁就写出了被引用超过十万次的论文（Vaswani、Gomez），
有些人从国际象棋神童成长为诺贝尔奖得主（Hassabis），还有些人放弃了
数十亿美元的商业机会转而投身AI安全研究（Sutskever、Bengio）。
他们的共同点是：每个人都在某个关键时刻做出了改变AI历史轨迹的贡献。

本章按主题分组介绍约三十位深度学习先驱，从``图灵奖三巨头''开始，
经过Transformer的八位作者，到更广泛的奠基性人物。
对于每一位，我们不仅讲述``做了什么''，更试图回答``为什么是他/她''——
什么样的背景、训练和直觉，使得这些人能够看到同时代人看不到的东西。


% =============================================================
\section{The Turing Award Trio: Hinton, LeCun, Bengio}
\label{sec:turing-trio}
% =============================================================

2018年3月27日，ACM（美国计算机协会）宣布将当年的图灵奖授予
Geoffrey Hinton、Yann LeCun和Yoshua Bengio，以表彰他们
``在深度神经网络方面做出的概念性和工程性突破''。
这是计算机科学的最高荣誉，被誉为``计算领域的诺贝尔奖''，
奖金100万美元。三位获奖者随后被媒体冠以``AI教父''
（Godfathers of AI）的称号——一个他们本人态度各异的标签。

\begin{whybox}[Why This Matters: 三十年的坚持]
在深度学习于2012年``爆发''之前，神经网络研究在AI主流社区中
长期处于边缘地位。从1980年代末到2000年代中期，支持向量机（SVM）、
核方法和概率图模型才是学术会议的主角。Hinton、LeCun和Bengio
在这段``AI寒冬''中坚持了近三十年，不断改进训练算法、
提出新的网络架构、培养学生，最终等到了算力和数据
追上他们理论的时刻。他们的故事是科学史上
``少数派报告最终成为主流共识''的经典案例。
\end{whybox}


% -------------------------------------------------------------
\subsection{Geoffrey Hinton——敲响警钟的教父}
\label{sec:hinton}
% -------------------------------------------------------------

\subsubsection{从心理学到AI}
Geoffrey Everest Hinton于1947年12月6日出生在英国伦敦温布尔登。
他的家族有着深厚的学术传统——曾曾祖父是逻辑学家George Boole
（布尔代数的发明者），曾曾祖母是数学教育家Mary Everest Boole。
他的中间名``Everest''来自另一位著名亲戚——
George Everest爵士，英属印度的测量总监，
珠穆朗玛峰的英文名``Mount Everest''正是以他命名。
父亲Howard Hinton是著名的昆虫学家，
叔叔Colin Clark是经济学家，
外曾祖父James Hinton则是外科医生和作家。
可以说，Boole-Hinton家族是一条贯穿两个世纪的学术谱系，
从布尔代数到深度学习，从逻辑的数学化到智能的计算化，
形成了一条惊人的思想传承链。

Hinton早年在布里斯托尔的Clifton College就读，
随后进入剑桥大学国王学院，于1970年获得实验心理学学士学位。

``我一直想理解大脑是如何工作的，''Hinton在多次采访中回忆道。
他在剑桥期间曾短暂考虑过木匠和建筑学，但最终被一个问题
深深吸引：\textbf{人类如何从经验中学习？}
这个问题引导他走向人工智能。
1978年，他在爱丁堡大学获得人工智能博士学位，
导师是Christopher Longuet-Higgins。

\subsubsection{与慢性背痛共处的人生}
Hinton的个人生活中有一个鲜为人知但深刻影响了他工作方式的细节：
他从青少年时期就患有严重的慢性背痛。
``我最后一次坐下是在2005年，''他经常这样说，
``而那是一个错误。''到了五十多岁，背痛严重到
他决定彻底放弃坐姿。此后他所有的工作都是站立完成的，
乘车时他横躺在后座上，吃饭时像僧侣一样跪在泡沫垫子上。
``如果你让它完全控制你的生活，它反而不会给你任何麻烦，''
他对《纽约时报》记者Cade Metz这样解释。

这个细节揭示了Hinton性格中的一个核心特质：
\textbf{他愿意彻底重构自己的生活方式来适应约束条件}——
无论是身体上的约束，还是智识上的约束。
正如他用三十年坚持神经网络研究一样，
他用二十年适应了一种大多数人无法想象的日常生活。

\subsubsection{反向传播与Boltzmann机}
博士毕业后，Hinton先后在苏塞克斯大学、加州大学圣迭戈分校和
卡内基梅隆大学工作。在这些年里，他做出了两项奠基性贡献：

\begin{itemize}[noitemsep]
  \item \textbf{反向传播的推广}（1986年）：与David Rumelhart和
    Ronald Williams合作发表的论文
    ``Learning Representations by Back-propagating Errors''
    发表在\textit{Nature}上，展示了如何通过梯度下降有效地训练
    多层神经网络。虽然反向传播的数学原理并非他们首创
    （可追溯至1960年代的控制论），但这篇论文首次清晰地
    证明了它在学习内部表示方面的强大能力。
  \item \textbf{Boltzmann机}（1985年）：与Terry Sejnowski合作
    发明的Boltzmann机是一种随机递归神经网络，
    基于统计物理中的能量函数和模拟退火。
    这项工作将物理学的工具引入了神经网络训练，
    也正是2024年诺贝尔物理学奖表彰的核心贡献之一。
\end{itemize}

\subsubsection{多伦多大学与``AI寒冬''的坚守}
1987年，由于对里根政府军事化AI研究的不满
（Hinton是坚定的和平主义者），他离开美国，
加入了加拿大高等研究院（CIFAR）并成为多伦多大学计算机科学系教授。
在随后的二十多年里，多伦多大学成为全球神经网络研究的圣地。

这段时期的关键里程碑包括：
\begin{itemize}[noitemsep]
  \item \textbf{深度信念网络}（2006年）：Hinton提出了一种
    逐层贪婪预训练的方法，首次证明深层网络可以被有效训练。
    这篇发表在\textit{Neural Computation}上的论文
    被广泛认为是深度学习复兴的起点。
  \item \textbf{AlexNet}（2012年）：Hinton的两位学生
    Alex Krizhevsky和Ilya Sutskever开发的AlexNet
    在ImageNet大规模视觉识别挑战赛中以压倒性优势获胜，
    top-5错误率从26.2\%骤降至15.3\%。
    这一结果如同一声惊雷，宣告了深度学习时代的到来。
\end{itemize}

\subsubsection{Google岁月与离开}
2013年，Google以约4400万美元收购了Hinton与学生创办的
DNNresearch Inc.。Hinton加入Google Brain，担任杰出研究员，
后升任副总裁和工程院士（Vice President \& Engineering Fellow）。
在Google期间，他继续推动深度学习研究，
同时保持在多伦多大学的职位。

然而，2023年5月，Hinton做出了一个震惊AI界的决定：
他从Google辞职，公开表达对AI安全的深切担忧。
``我对自己毕生的工作感到一些后悔，''他对《纽约时报》说。
``我用一个常见的借口来安慰自己：如果我不做，
别人也会做。''他警告说，AI系统可能很快就会比人类更聪明，
而我们还没有准备好应对这种情况。

\subsubsection{双料诺贝尔级荣誉}
2024年10月8日，瑞典皇家科学院宣布将2024年诺贝尔物理学奖
授予Geoffrey Hinton和John Hopfield，
``以表彰他们在利用人工神经网络实现机器学习方面的
基础性发现和发明''。Hinton的Boltzmann机——
建立在Hopfield网络基础上的随机生成模型——
被认为是现代深度学习的物理学根基。

这使得Hinton成为极少数同时获得图灵奖和诺贝尔奖的科学家之一，
也是首位因AI研究获得诺贝尔物理学奖的人。

得知获奖的那一刻颇具戏剧性。Hinton当时正在加州一家
``没有网络连接、手机信号也很差''的廉价旅馆里睡觉，
一个带有瑞典口音的国际电话把他叫醒。
``有人问我是不是Geoffrey Hinton……
然后告诉我获得了诺贝尔物理学奖。
我的第一反应是：等一下，我不做物理学。''
此后好几天，他都在用``统计推理''计算
一个心理学家获得诺贝尔物理学奖的概率有多低，
以此验证这件事是否真的发生了。
他原本当天要去做核磁共振检查，
``我想我得取消了，''他说。

\subsubsection{AI安全的具体警告}
Hinton对AI安全的担忧不是笼统的恐惧，而是指向了多个具体风险。
在2024年12月10日的诺贝尔奖晚宴致辞中，他系统地阐述了这些担忧：

\begin{itemize}[noitemsep]
  \item \textbf{短期风险}：AI已经通过推送令人愤怒的内容
    创造了分裂性的信息茧房；它正被威权政府用于大规模监控，
    被网络犯罪分子用于钓鱼攻击。
  \item \textbf{中期风险}：AI可能很快被用于制造新型病毒，
    或开发能\textbf{自主选择目标}的自主武器。
  \item \textbf{长期风险}：人类正走在创造出
    ``比我们更聪明的数字存在''的道路上，
    而我们完全不知道是否能控制它们。
    如果这些系统由\textbf{追求短期利润}的公司创造，
    公共安全将不会是首要考量。
\end{itemize}

``所有这些短期风险都需要各国政府和国际组织
紧急而有力的关注，''他在致辞中说。
这使得他的诺贝尔演说成为了一次罕见的``获奖者警告世界''时刻——
在庆典的华丽氛围中，他选择用最直白的语言说出最沉重的话。

\begin{corebox}[Core Insight: Hinton的遗产]
Hinton的职业生涯完美诠释了一个悖论：
\textbf{最深刻地推动了一项技术发展的人，也可能最深刻地理解
其危险性。}他的故事提醒我们，科学家对自己创造物的责任
不会因为技术的成功而消失——恰恰相反，成功使得这种责任
变得更加沉重。截至2026年，Hinton继续担任
Vector Institute的首席科学顾问，
同时积极参与全球AI安全政策讨论。
\end{corebox}

\bigskip

% -------------------------------------------------------------
\subsection{Yann LeCun——从手写识别到世界模型}
\label{sec:lecun}
% -------------------------------------------------------------

\subsubsection{法国少年的神经网络梦}
Yann Andr\'{e} LeCun（法语发音接近``勒坤''，
但他本人接受英语发音``le-KUNN''）于1960年7月8日出生在
法国巴黎北郊的Soisy-sous-Montmorency。
他在ESIEE Paris获得电气工程师文凭（1983年），
然后在巴黎第六大学（Pierre and Marie Curie University，
现索邦大学）攻读博士学位，1987年毕业。
他的博士论文提出了反向传播算法的早期变体，
与Hinton的工作平行但独立发展。

博士毕业后，LeCun前往多伦多大学跟随Hinton做博士后研究，
这段经历奠定了两人终身的学术友谊。

\subsubsection{AT\&T Bell Labs与卷积神经网络}
1988年，LeCun加入AT\&T Bell Labs，
开始了他最具影响力的一段研究生涯。在这里，
他开发了\textbf{卷积神经网络}（CNN）的核心架构，
特别是用于手写数字识别的LeNet系列。
LeNet-5（1998年）能够准确识别手写支票上的数字，
被美国银行系统广泛部署——在AI的``寒冬''时期，
这是少数几个真正商业化的神经网络应用之一。

LeCun在Bell Labs期间的其他重要贡献包括：
\begin{itemize}[noitemsep]
  \item DjVu图像压缩技术（与L\'{e}on Bottou合作）
  \item Optimal Brain Damage——一种网络剪枝方法
  \item Graph Transformer Networks——将CNN推广到图结构
\end{itemize}

\subsubsection{纽约大学与Meta AI}
2003年，LeCun加入纽约大学（NYU）Courant数学科学研究所
担任教授。他后来成为NYU数据科学中心的创始主任，
目前持有Jacob T. Schwartz计算机科学讲席教授头衔。

2013年，LeCun被Facebook（现Meta）聘为AI研究部门
（Facebook AI Research，FAIR）的首任主任。
2018年起，他担任Meta的副总裁兼首席AI科学家
（VP \& Chief AI Scientist）。在Meta期间，
他主导了开源大语言模型LLaMA系列的研发方向，
推动了``开放AI''的理念。

值得一提的是，LeCun是AI领域最活跃的``社交媒体战士''之一。
他在X（原Twitter）上以犀利、直率甚至好斗的风格著称，
经常与AI安全倡导者（包括他的图灵奖共同获奖者Hinton和Bengio）
公开辩论。他坚持认为当前的大语言模型
``在本质上不可能达到人类水平的智能''，
因为它们缺乏对物理世界的理解。

\subsubsection{JEPA：一个与LLM截然不同的技术路线}
LeCun的技术愿景集中在\textbf{JEPA}
（Joint Embedding Predictive Architecture，联合嵌入预测架构）上。
与大语言模型通过预测下一个token来学习不同，
JEPA在\textbf{抽象表示空间}中进行预测——
它不试图重建输入信号的每一个细节，
而是学习预测事物的抽象表征。

LeCun用一个生动的比喻来解释这个区别：
``一只猫理解物理世界的能力比任何大语言模型都强。
猫知道如果一个物体从桌子边缘落下，它会掉到地上。
但没有一个LLM真正`理解'这一点——
它们只是在统计上模仿了描述这件事的文本。''

他的论点是：对于文本这种低带宽、高信息密度的媒介，
逐token预测是有效的。但对于视频和真实世界数据
这种高带宽、高冗余的信号，生成式方法
（如自编码器和掩码图像建模）会系统性地失败，
因为细节的不确定性太大了。
JEPA通过在嵌入空间而非像素空间中预测，
优雅地绕过了这个问题。
2024年发布的V-JEPA 2已在物理推理基准测试中
达到了最先进的水平。

\subsubsection{``AI辩论场上的格斗家''}
LeCun在社交媒体上的辩论风格在AI社区中是独一无二的。
他不回避与任何人——包括他的图灵奖共获者——进行
长篇、公开、有时近乎尖锐的学术辩论。
当Hinton警告``AI可能很快比人类更聪明''时，
LeCun的回应通常是：``我们连猫的智能水平都还没达到。''
当AI末日论者（doomers）呼吁暂停AI研究时，
他称这种立场``基于对技术现状的根本误解''。

这种风格让他赢得了一批铁杆拥趸，也招致了不少批评。
但LeCun似乎享受这种辩论——他将自己定位为
AI社区中``理性乐观主义''的代言人，
认为恐惧和过度监管比AI本身更危险。

\subsubsection{离开Meta，创办AMI Labs}
LeCun离开Meta的背景比简单的``创业冲动''更复杂。
2025年6月，Mark Zuckerberg投入143亿美元收购Scale AI，
并在Meta内部成立了以LLM为核心的
Meta Superintelligence Labs新部门。
LeCun一手创建的FAIR（Facebook AI Research）被边缘化，
机器人团队被解散。到2025年11月，
LeCun的世界观与Meta的技术方向之间的裂痕已无法弥合。

``我走进Mark的办公室，告诉他我要离开了，''
LeCun后来回忆这次对话。
他随即创立了\textbf{AMI Labs}
（Advanced Machine Intelligence Labs），
总部设在巴黎，CEO是Nabla创始人Alexandre LeBrun。
AMI Labs的目标是开发能够\textbf{理解物理世界、
具有持久记忆、能够推理并规划复杂行动}的AI系统——
即他长期倡导的``世界模型''（World Models）。

2026年3月，AMI Labs宣布完成10.3亿美元融资，
投前估值35亿美元，由Cathay Innovation、
Greycroft、Bezos Expeditions等共同领投。
``我的预测是，`世界模型'将成为下一个热词，''
LeBrun对TechCrunch说。
``六个月内，每家公司都会自称是世界模型公司来融资。''
这标志着65岁的LeCun开启了职业生涯的新篇章——
他正在用自己的声誉和智识资本，
押注整个AI产业的方向是错误的。

\begin{keybox}[Key Awards: Yann LeCun]
\begin{itemize}[noitemsep]
  \item ACM图灵奖（2018年，与Hinton、Bengio共享）
  \item IEEE神经网络先驱奖（2014年）
  \item 美国国家工程院院士（2017年）
  \item 美国国家科学院院士（2021年）
  \item 法国荣誉军团勋章（2023年）
  \item Queen Elizabeth Prize for Engineering（2025年）
\end{itemize}
\end{keybox}

\bigskip

% -------------------------------------------------------------
\subsection{Yoshua Bengio——从深度学习到AI安全}
\label{sec:bengio}
% -------------------------------------------------------------

\subsubsection{巴黎出生，蒙特利尔成长}
Yoshua Bengio于1964年3月5日出生在法国巴黎，
随家人移居加拿大。他在麦吉尔大学完成了全部学位：
计算机工程学士（1986年）、计算机科学硕士（1988年）
和计算机科学博士（1991年）。
博士毕业后，他先后在MIT脑与认知科学系和AT\&T Bell Labs
做博士后研究，然后于1993年加入蒙特利尔大学计算机科学系，
此后再未离开。

Bengio的弟弟Samy Bengio同样是杰出的机器学习研究者
（曾在Google工作，现在Apple），兄弟二人的论文经常被混淆。

\subsubsection{语言模型与注意力机制的奠基者}
Bengio的研究贡献涵盖深度学习的多个核心领域：

\begin{itemize}[noitemsep]
  \item \textbf{神经网络语言模型}（2003年）：论文
    ``A Neural Probabilistic Language Model''首次提出用神经网络
    学习词的分布式表示（word embeddings）并建模语言概率，
    直接启发了后来的Word2Vec和所有现代大语言模型。
  \item \textbf{注意力机制}（2014年）：与Dzmitry Bahdanau合作的
    论文``Neural Machine Translation by Jointly Learning to Align
    and Translate''引入了注意力机制（Attention），
    允许模型在翻译时动态地``关注''输入序列的不同部分。
    这一思想后来被Transformer架构推广为``自注意力''（Self-Attention），
    成为现代AI的核心计算原语。
  \item \textbf{生成对抗网络}（2014年）：Bengio作为Ian Goodfellow
    的博士导师，参与了GAN论文的指导。
  \item \textbf{去噪自编码器}（2008年）：为后来的扩散模型
    和自监督学习奠定了基础。
\end{itemize}

\subsubsection{Mila与AI生态系统建设}
2017年，Bengio创立了Mila——魁北克人工智能研究所，
并担任科学主任直到2025年3月转为科学顾问。
Mila已发展成为全球最大的深度学习学术研究中心之一，
拥有超过1000名研究人员。

Bengio同时是CIFAR``机器与大脑学习''项目的联合主任，
IVADO的创始科学主任，
以及加拿大CIFAR AI讲席教授。
他的影响力不仅在于自身研究，更在于培养了一代又一代
深度学习研究者。

\subsubsection{全球最高引学者与AI安全转向}
截至2025年11月，Bengio成为首位Google Scholar引用量
突破100万的AI研究者，也是全领域在世学者中引用量最高的人。
他的h-index和总引用量均居全球计算机科学家之首。

然而，大约从2023年开始，Bengio将越来越多的精力
投入到AI安全和治理领域。他签署了多封公开信呼吁
暂停大型AI实验，在联合国、G7和各国政府面前作证，
并于2024年被\textit{TIME}杂志评为全球100位最具影响力人物之一。

2025年6月，Bengio宣布创立\textbf{LawZero}——
一个致力于``安全优先设计''（safe-by-design）AI的非营利组织。
LawZero的成立源于他对当前前沿AI模型日益增长的危险能力的担忧，
包括欺骗、自我保存和目标失调等行为。Bengio强调，
LawZero作为非营利组织的结构是为了``隔绝市场和政府压力，
避免AI安全被商业利益所妥协''。

\begin{whybox}[Why This Matters: 三位教父的分歧]
三位图灵奖共同获得者在AI安全问题上的巨大分歧，
是理解当前AI辩论的一把钥匙。Hinton离开Google，
公开表达``生存风险''级别的担忧；Bengio创立LawZero，
致力于技术层面的安全解决方案；而LeCun则坚持认为
当前的AI系统离人类水平智能还很远，过度恐慌反而有害。
这三种立场——\textbf{警报者}、\textbf{建设性安全主义者}
和\textbf{务实乐观主义者}——代表了AI社区内部最重要的
三种思潮。读者在阅读本书姊妹篇\textit{Emergence}
第十三章（前沿问题）时，可以进一步了解这场辩论的技术细节。
\end{whybox}


% =============================================================
\section{``Attention Is All You Need''：八位Transformer作者}
\label{sec:transformer-authors}
% =============================================================

2017年6月12日，一篇由八位Google研究员撰写的论文
被提交到arXiv，标题简洁而大胆：``Attention Is All You Need''。
论文提出了\textbf{Transformer}架构——一种完全基于
自注意力机制（Self-Attention）的序列建模方法，
摒弃了此前主导序列任务的循环神经网络（RNN）和
长短期记忆网络（LSTM）。

这篇论文截至2026年已被引用超过15万次，
成为人类历史上被引用最多的计算机科学论文之一。
GPT、BERT、PaLM、LLaMA、Claude、Gemini、DALL-E、
Stable Diffusion、AlphaFold 2、Codex、Whisper——
现代AI工业的整个产品线，都运行在这八个人发明的架构之上。

更引人注目的是：到2023年8月，八位作者已全部离开Google，
分散在各自创立或加入的创业公司中。他们的离去象征着一个时代的结束
——Google曾经是AI人才的最大聚集地，
但它未能将这项革命性发明转化为先发优势。

\begin{corebox}[Core Insight: 一篇论文改变世界]
``Attention Is All You Need''的影响力可以与计算机科学史上
少数几篇论文相提并论：图灵1936年关于可计算数的论文、
香农1948年关于信息论的论文、以及Backus等人1957年的
Fortran报告。它不仅解决了一个技术问题（高效的序列建模），
更开启了一个全新的计算范式。
关于Transformer架构的技术细节，
读者可参阅姊妹篇\textit{Emergence}第二章的深入讨论。
\end{corebox}


% -------------------------------------------------------------
\subsection{Ashish Vaswani——第一作者与Essential AI}
\label{sec:vaswani}
% -------------------------------------------------------------

Ashish Vaswani是``Attention Is All You Need''的第一作者，
也是Transformer架构的核心设计者。他在印度长大，
在印度理工学院（IIT）完成本科后赴美深造，
在南加州大学（USC）获得计算机科学博士学位，
专攻机器学习与自然语言处理。

加入Google Brain后，Vaswani花了数年时间研究注意力机制
在神经网络中的应用。他意识到一个关键问题：
RNN和LSTM的\textbf{顺序计算}本质严重限制了训练效率——
每个时间步必须等待前一步完成，无法充分利用GPU的并行能力。
Transformer的核心创新正是用自注意力完全取代循环计算，
使得序列中所有位置可以同时交互。

论文发表于NeurIPS 2017（当时仍称NIPS）后，Vaswani继续在Google
工作了一段时间，随后于2022年与Niki Parmar共同创立了
\textbf{Essential AI}，一家面向企业的AI公司。
Essential AI在2024年获得了约5600万美元的融资。
2025年12月，Essential AI发布了首个模型Rnj-1
（以印度数学家Ramanujan命名），
号称仅用80亿参数就达到了前沿水平的编码、数学和
智能体推理能力——远小于OpenAI的万亿参数量级。

\bigskip

% -------------------------------------------------------------
\subsection{Noam Shazeer——数学天才与Character.AI传奇}
\label{sec:shazeer}
% -------------------------------------------------------------

Noam Shazeer或许是八位作者中最具传奇色彩的人物。
出生于约1975--1976年，他在1994年以满分成绩
获得国际数学奥林匹克（IMO）金牌——
那一年美国队实现了罕见的全队满分壮举，
这一纪录直到2022年才被中国队追平。
\textit{TIME}杂志曾报道这一成就。

Shazeer在杜克大学（Duke University）学习计算机科学和数学，
2000年毕业后加入了成立仅两年的Google，
成为最早的几百名员工之一。在Google的二十年间，
他从软件工程师做到高级软件工程师和研究科学家，
参与了Google搜索引擎的核心算法改进，
是Google广告拼写校正系统和早期机器翻译系统的关键贡献者。

2017年，当``Attention Is All You Need''的早期草稿出来时，
Shazeer惊讶地发现自己的名字被列在第一位
（最终发表时按字母排序）。``我当时并没有想太多，''他回忆说。
但历史证明，他对Transformer的贡献是不可或缺的。

2021年，Shazeer离开Google，与Daniel De Freitas共同创立了
\textbf{Character.AI}——一个允许用户与AI角色对话的平台。
Character.AI迅速走红，一度成为增长最快的消费级AI应用之一。
然而，2024年8月，在一项惊人的交易中，
Google以约27亿美元的价格将Shazeer``请回''：
Google获得了Character.AI的技术授权和Shazeer本人，
而Character.AI则保持独立运营。
Shazeer重返Google DeepMind担任副总裁和AI研究负责人，
据报道目标直指AGI——通用人工智能。

\bigskip

% -------------------------------------------------------------
\subsection{其他六位作者的去向}
\label{sec:other-transformer-authors}
% -------------------------------------------------------------

\subsubsection{Llion Jones与Sakana AI}
Llion Jones是一位威尔士裔英国人，在伯明翰大学获得人工智能与
计算机科学学士学位及高级计算机科学硕士学位。
他先在YouTube工作了三年多，2015年加入Google，
专注于机器智能和自然语言理解的研究。

Jones在Google工作了近12年，是八位作者中最后一位离开的。
2023年8月，他远赴日本东京，
与前Google研究科学家David Ha共同创立了
\textbf{Sakana AI}（``sakana''在日语中意为``鱼''）。
Sakana AI以``自然启发的AI研究''为理念，
探索超越Transformer的下一代AI架构。
公司在2024年获得了超过2亿美元的融资。

\subsubsection{Jakob Uszkoreit与Inceptive}
Jakob Uszkoreit出生于德国，在柏林工业大学获得硕士学位。
他是八位作者中唯一一位将AI转向生命科学领域的人。
在Google工作13年后（2008--2021），
他先后参与了Google Translate的早期开发、
Google Assistant的语言理解团队建设以及Google Brain的深度学习研究。

2021年7月，Uszkoreit与斯坦福大学生物化学家Rhiju Das
共同创立了\textbf{Inceptive}，
一家用深度学习和高通量实验``学习生命语言''的生物技术公司。
Inceptive专注于RNA药物设计——
利用类似Transformer的架构来理解和设计RNA分子。
``如果Transformer可以理解人类语言，''Uszkoreit说，
``它也应该能理解生命的语言。''

值得一提的是，Uszkoreit的父亲Hans-J\"{u}rgen Uszkoreit
是计算语言学领域的知名教授，可以说NLP研究在这个家族中``遗传''了。

\subsubsection{Aidan Gomez与Cohere}
Aidan Gomez是八位作者中最年轻的，论文发表时年仅20岁。
1998年（一说1997年）出生于加拿大多伦多，
他当时还是多伦多大学计算机科学本科生，
在Google Brain多伦多办公室实习时参与了Transformer论文的研究。
他曾与Geoffrey Hinton和Jeff Dean等AI传奇人物共事。

实习结束后，Gomez完成了本科学业，
并开始在牛津大学攻读计算机科学博士学位。
2019年，他与Ivan Zhang和Nick Frosst（Hinton的学生）
共同创立了\textbf{Cohere}——一家面向企业的AI语言模型公司。
截至2024年，Cohere的估值已达约55亿美元，
由Nvidia等顶级投资者支持，
成为加拿大最成功的AI创业公司之一。

\subsubsection{Illia Polosukhin与NEAR Protocol}
Illia Polosukhin来自乌克兰，在哈尔科夫国立无线电电子大学
获得应用数学与计算机科学硕士学位。
他在Google Research工作了约三年（2014--2017），
参与了TensorFlow的开发，并领导了Google搜索引擎
问答功能的团队。

Polosukhin是八位作者中唯一一位将事业转向
区块链和Web3的人。2017年，他与Alexander Skidanov
创立了NEAR.ai，最初专注于AI和程序合成。
但很快，两人看到了去中心化技术的潜力，
将项目转型为\textbf{NEAR Protocol}——
一个Layer-1区块链平台，于2020年正式上线。
到2025年，NEAR已成为Web3领域最重要的基础设施项目之一。

\subsubsection{Niki Parmar——从自学编程到Anthropic}
Niki Parmar的故事是八位作者中最励志的之一。
她是印度裔，没有考上印度理工学院（IIT），
一度认为自己``错过了机会''。但她选择了自学编程，
在南加州大学（USC）获得计算机科学硕士学位后，
于2015年24岁时加入Google——成为团队中最年轻的成员，
也是唯一一个没有博士学位的人。

在Google Brain，Parmar成为Transformer论文的核心贡献者之一，
并领导了将Transformer扩展到大规模语言模型的研究。
离开Google后，她先后共同创立了两家AI公司：
\textbf{Adept AI}（2021年，获得超过4亿美元融资）和
\textbf{Essential AI}（2023年，与Vaswani共同创立）。
2025年1月，她加入了\textbf{Anthropic}
担任技术人员（Member of Technical Staff），
为Claude系列模型的开发做出贡献。

\subsubsection{\L{}ukasz Kaiser——从数学逻辑到推理模型}
\L{}ukasz Kaiser于1981年12月24日出生在波兰弗罗茨瓦夫。
他在弗罗茨瓦夫大学学习计算机科学和数学，
然后在德国亚琛工业大学（RWTH Aachen）获得博士学位，
研究方向是数学逻辑、自动机理论和自动结构——
这些是理论计算机科学的核心领域。

博士毕业后，Kaiser在巴黎第七大学（Paris Diderot）担任
终身研究员，2013年转入Google Brain，
在那里他共同设计了用于机器翻译、解析和其他生成任务的
神经模型，并参与了TensorFlow系统和Tensor2Tensor库的开发。

Kaiser后来加入了\textbf{OpenAI}，
在那里他参与了ChatGPT、GPT-4以及推理模型o1和o3的开发。
在2025年TEDAI大会上，他被介绍为
``共同发明了推理模型（o1、o3）、ChatGPT和GPT-4，
并共同开创了Transformer和TensorFlow系统''的人。

\begin{keybox}[Key Reference: 八位作者的去向与公司估值（截至2026年初）]
\begin{itemize}[noitemsep]
  \item \textbf{Ashish Vaswani} $\to$ Essential AI
    （联合创始人兼CEO，估值数亿美元级）
  \item \textbf{Noam Shazeer} $\to$ Character.AI $\to$
    回归Google DeepMind（VP，Character.AI交易价值约27亿美元）
  \item \textbf{Niki Parmar} $\to$ Adept AI $\to$ Essential AI $\to$
    Anthropic（Member of Technical Staff）
  \item \textbf{Jakob Uszkoreit} $\to$ Inceptive
    （联合创始人，RNA药物设计）
  \item \textbf{Llion Jones} $\to$ Sakana AI
    （联合创始人，东京，融资超2亿美元）
  \item \textbf{Aidan Gomez} $\to$ Cohere
    （联合创始人兼CEO，估值约55亿美元）
  \item \textbf{Illia Polosukhin} $\to$ NEAR Protocol
    （联合创始人，Layer-1区块链）
  \item \textbf{\L{}ukasz Kaiser} $\to$ OpenAI
    （参与GPT-4、ChatGPT和o1/o3推理模型开发）
\end{itemize}
八位作者创立或参与的公司总估值超过\textbf{百亿美元}——
而这一切都起源于2017年的一篇论文。
\end{keybox}


% =============================================================
\section{LSTM双雄：Schmidhuber与Hochreiter}
\label{sec:lstm-duo}
% =============================================================

在Transformer统治AI世界之前，有一种架构在序列建模领域
称霸了近二十年——\textbf{长短期记忆网络}
（Long Short-Term Memory，LSTM）。
LSTM的两位发明者——J\"{u}rgen Schmidhuber和Sepp Hochreiter
——的故事既是合作的佳话，也是AI学术界优先权之争的缩影。


% -------------------------------------------------------------
\subsection{J\"{u}rgen Schmidhuber——``如果AI成熟了，它会叫我爸爸''}
\label{sec:schmidhuber}
% -------------------------------------------------------------

J\"{u}rgen Schmidhuber于1963年1月17日出生在德国慕尼黑。
他在慕尼黑工业大学（TU Munich）完成了全部学业。
从15岁起，他就有一个明确的人生目标：
\textbf{``建造一个比我更聪明的自我改进AI，然后退休。''}
这句话成了他数十年来的签名式宣言。

Schmidhuber目前担任瑞士Dalle Molle人工智能研究所（IDSIA）
的科学主任，以及沙特阿拉伯KAUST（阿卜杜拉国王科技大学）
AI Initiative的主任和计算机科学教授。

\subsubsection{``奇迹年''1990--1991}
Schmidhuber将1990--1991年称为他的``Annus Mirabilis''（奇迹年），
声称在这两年里奠定了现代深度学习的多项基础：
\begin{itemize}[noitemsep]
  \item \textbf{线性Transformer}的原理（1991年），
    他认为这是自注意力机制的前身
  \item \textbf{生成对抗思想}（1990年），
    通过对抗式神经网络实现人工好奇心
  \item 神经网络蒸馏、深度学习、元学习等概念的早期形式
\end{itemize}

这些优先权声明在AI社区中引起了持续的争议。
Schmidhuber以``频繁且强烈地主张自己的优先权''著称，
曾公开批评Transformer论文的作者未充分引用他的早期工作，
也曾与GAN的发明者Ian Goodfellow发生过公开争论。
《纽约时报》为他写过一篇著名的特稿，标题是：
``When A.I. Matures, It May Call J\"{u}rgen Schmidhuber `Dad'\,''。

\subsubsection{LSTM与实际影响}
不管优先权之争如何，一个不争的事实是：
LSTM（1997年论文，Hochreiter为第一作者，Schmidhuber为通讯作者）
在2010年代成为了自然语言处理和语音识别的主导技术。
到2017年，LSTM运行在超过30亿台智能手机上：
Google的语音识别（基于CTC-LSTM）、Google翻译（每天超过40亿次
LSTM翻译）、Apple的Siri和QuickType、Amazon的Alexa
——全部依赖于LSTM架构。

Schmidhuber的实验室还在2011年首次以深度神经网络赢得了
官方计算机视觉竞赛（达到超人表现），
在2012年赢得了首个深度学习医学影像竞赛（癌症检测）。
这些突破比AlexNet还要早几个月，
但由于AlexNet在ImageNet上的戏剧性胜利获得了更多关注，
Schmidhuber的贡献常被忽视——这也是他频繁发声的原因之一。

\bigskip

% -------------------------------------------------------------
\subsection{Sepp Hochreiter——梯度消失问题的发现者}
\label{sec:hochreiter}
% -------------------------------------------------------------

Josef ``Sepp'' Hochreiter于1967年2月14日出生在德国M\"{u}hldorf。
他在慕尼黑工业大学学习，在Schmidhuber的指导下
完成了改变深度学习历史的毕业论文。

\subsubsection{发现梯度消失问题}
1991年，年轻的Hochreiter在他的毕业论文中
\textbf{数学地证明了}标准循环神经网络的致命缺陷：
梯度在反向传播过程中会指数级地衰减
（vanishing gradient problem），使得网络
无法学习超过约10个时间步的长距离依赖关系。
这不是一个小麻烦——它是一个根本性的障碍。

四年后，1997年，Hochreiter和Schmidhuber发表了他们的解决方案：
\textbf{LSTM}。LSTM的核心创新是引入了``门控''机制
（输入门、遗忘门和输出门），允许网络选择性地
记住或遗忘信息，从而在长序列中保持梯度流动。
这篇发表在\textit{Neural Computation}上的论文
成为了AI历史上被引用最多的论文之一。

\subsubsection{学术生涯与xLSTM}
Hochreiter后来在柏林工业大学和科罗拉多大学博尔德分校工作过，
1999年在慕尼黑工业大学获得博士学位（导师为Wilfried Brauer）。
2006年起，他在奥地利约翰内斯开普勒大学林茨分校（JKU Linz）任教，
先后领导了生物信息学研究所和机器学习研究所，
并创立了Linz Institute of Technology AI Lab。
他也是人工智能高等研究所（IARAI）的创始主任。

2024年，Hochreiter提出了\textbf{xLSTM}
（Extended LSTM），声称这种改进版的LSTM在某些任务上
可以与Transformer竞争甚至超越它们。
``我们还处于早期阶段，''他说，
``LSTM的故事还没有结束。''


% =============================================================
\section{ImageNet时代的英雄}
\label{sec:imagenet-heroes}
% =============================================================

如果说LSTM统治了序列建模的世界，
那么\textbf{卷积神经网络}（CNN）则征服了计算机视觉。
而这场征服的导火索，是2012年的ImageNet竞赛。

% -------------------------------------------------------------
\subsection{Alex Krizhevsky——AlexNet与深度学习革命的导火索}
\label{sec:krizhevsky}
% -------------------------------------------------------------

Alex Krizhevsky是一位加拿大计算机科学家，
出生在乌克兰，在多伦多大学Hinton的实验室攻读博士学位。
2012年，他开发了一个被命名为\textbf{AlexNet}的
深层卷积神经网络，并用两块GeForce GTX 580显卡进行训练。

AlexNet在ImageNet大规模视觉识别挑战赛上的表现
是计算机视觉历史上最戏剧性的时刻之一：
top-5错误率15.3\%，而第二名使用传统手工特征的方法
只达到26.2\%——差距将近11个百分点。
这不像是渐进式的改进，而像是一种完全不同的技术。

``Ilya觉得我们应该做这件事，Alex让它跑起来了，
而我得了诺贝尔奖，''Hinton后来开玩笑说。
确实，是Sutskever的直觉——
``神经网络的性能会随数据量扩展''——
和Krizhevsky的工程才华的结合，
创造了这个历史性的突破。

AlexNet之后不久，Krizhevsky和Sutskever将他们的公司
DNN Research Inc.卖给了Google。Krizhevsky在Google工作到
2017年9月，因``对工作失去兴趣''而离开，
转入了一家名为Dessa的小型深度学习公司。
他也是CIFAR-10和CIFAR-100数据集的主要创建者。
虽然他后来淡出了公众视野，但AlexNet的影响是永恒的：
它点燃了深度学习革命的导火索。

\bigskip

% -------------------------------------------------------------
\subsection{Fei-Fei Li——ImageNet的创造者与``AI教母''}
\label{sec:feifei-li}
% -------------------------------------------------------------

如果AlexNet是点燃革命的火花，
那么\textbf{ImageNet}就是等待引爆的火药库。
而创造这个火药库的人，是被称为``AI教母''
（Godmother of AI）的李飞飞（Fei-Fei Li）。

\subsubsection{从新泽西干洗店到斯坦福讲台}
李飞飞于1976年出生在中国，16岁时随父母移民美国新泽西州。
家境清贫，她十几岁时就在父母的干洗店里帮工。
凭借卓越的学术能力，她获得奖学金进入普林斯顿大学，
1999年以优异成绩（High Honors）获得物理学学士学位。
随后她进入加州理工学院（Caltech）攻读博士学位，
2005年获得电气工程博士学位。

\subsubsection{ImageNet——一个曾被嘲笑的想法}
2007年，已是普林斯顿助理教授的李飞飞产生了一个大胆的想法：
创建一个包含数百万张标注图像的超大规模数据库，
用于训练和评估计算机视觉算法。
当时的主流做法是在几千张图片的小数据集上做实验，
她的同事们认为这个想法``太野心勃勃了''。

2009年，她在迈阿密的一次学术会议上以海报形式
首次展示了ImageNet，几乎没有引起关注。
但她坚持了下来。ImageNet最终包含了约1500万张图片，
涵盖22000个类别，并从2010年开始举办年度竞赛
（ImageNet Large Scale Visual Recognition Challenge，ILSVRC）。

正是这个竞赛为AlexNet提供了舞台。
没有ImageNet的大规模数据，AlexNet的深度学习方法
就无法展示其压倒性的优势。在这个意义上，
\textbf{李飞飞的数据远见与Hinton的算法理论同样重要}。

\subsubsection{斯坦福HAI与World Labs}
李飞飞从2009年起在斯坦福大学任教，先后担任
斯坦福AI实验室（SAIL）主任（2013--2018年）。
2017--2018年间，她曾赴Google Cloud担任副总裁兼AI/ML首席科学家。
2019年，她共同创立了斯坦福以人为本AI研究院
（Stanford HAI），成为AI伦理和治理领域的重要声音。

2024年9月，李飞飞与三位同事创办了\textbf{World Labs}，
一家专注于``空间智能''（Spatial Intelligence）的AI公司，
获得了2.3亿美元融资（投资者包括Geoffrey Hinton），
估值达到约12.5亿美元。
2025年，她再次入选\textit{TIME} AI 100
最具影响力人物榜单。

\begin{keybox}[Key Lesson: 数据的力量]
李飞飞的ImageNet故事揭示了AI进步的一个常被忽视的维度：
\textbf{数据基础设施与算法创新同样重要}。
当学术界都在优化算法时，她选择了构建数据集——
一个看起来``不够性感''的工作。
但正是这种基础设施的远见，
为整个深度学习革命提供了必要条件。
\end{keybox}


% =============================================================
\section{AI大众化先锋：Andrew Ng与Andrej Karpathy}
\label{sec:popularizers}
% =============================================================

深度学习的成功不仅需要算法突破和工程创新，
还需要有人将这些知识传播给全世界。
Andrew Ng和Andrej Karpathy是AI教育和普及化的两面旗帜。

% -------------------------------------------------------------
\subsection{Andrew Ng——八百万学生的AI老师}
\label{sec:andrew-ng}
% -------------------------------------------------------------

吴恩达（Andrew Yan-Tak Ng）于1976年4月18日
出生在英国伦敦。他的父亲Ronald Paul Ng是血液学家和讲师。
吴恩达在卡内基梅隆大学获得计算机科学学士学位，
在MIT获得硕士学位，在加州大学伯克利分校获得博士学位。

\subsubsection{Google Brain与Coursera}
2011年，吴恩达在斯坦福大学开设了一门在线机器学习课程，
吸引了超过10万名学生注册——这在当时是前所未有的数字。
这一成功直接催生了\textbf{Coursera}（2012年，
与Daphne Koller共同创立），
如今是全球最大的MOOC（大规模开放在线课程）平台。

同时期，吴恩达创立并领导了\textbf{Google Brain}项目——
使用大规模深度学习算法的团队。
2012年，Google Brain产出了著名的``Google猫''实验：
一个拥有10亿参数的神经网络自主从YouTube视频中
学会了识别猫的概念，无需人工标注。
这一成果登上了各大媒体头条，
向公众展示了深度学习的无监督学习潜力。

\subsubsection{百度首席科学家}
2014年，吴恩达加入百度担任副总裁兼首席科学家，
建设了一支超过1300人的AI团队，
负责百度的AI技术和战略。2017年离开百度后，
他创立了多家公司：
\begin{itemize}[noitemsep]
  \item \textbf{DeepLearning.AI}：在线AI教育平台，
    至今已有超过800万学生学习过他的课程
  \item \textbf{Landing AI}：为制造业提供AI解决方案的SaaS公司
  \item \textbf{AI Fund}：一个初始规模1.75亿美元的
    AI创业投资基金
\end{itemize}

2024年4月，Amazon宣布任命吴恩达为其董事会成员，
标志着他在商业领域的影响力进一步扩大。

吴恩达被\textit{TIME}杂志评为2012年100位最具影响力人物之一，
并入选2023年TIME100 AI最具影响力AI人物榜单。
他目前是斯坦福大学兼职教授，
曾担任斯坦福AI实验室（SAIL）主任。

\begin{corebox}[Core Insight: 民主化的力量]
吴恩达的最大贡献或许不是任何一篇论文或一个产品，
而是\textbf{让全世界任何有互联网连接的人都能免费学习AI}。
从Google Brain到Coursera到DeepLearning.AI，
他始终致力于降低AI教育的门槛。在一个充斥着
``AI将取代人类''恐惧的时代，
吴恩达选择的回应是：``让更多人理解AI。''
\end{corebox}

\bigskip

% -------------------------------------------------------------
\subsection{Andrej Karpathy——从斯坦福到特斯拉再到教育}
\label{sec:karpathy}
% -------------------------------------------------------------

Andrej Karpathy于1986年10月23日出生在
捷克斯洛伐克布拉迪斯拉发（现斯洛伐克），
15岁时随家人移民到加拿大多伦多。
他在多伦多大学获得计算机科学和物理学双学位（2009年），
在英属哥伦比亚大学（UBC）获得硕士学位（2011年），
在斯坦福大学获得计算机科学博士学位（2016年），
导师正是李飞飞。

\subsubsection{OpenAI、Tesla、再回OpenAI}
Karpathy的职业轨迹是AI产业快速变化的缩影：
\begin{itemize}[noitemsep]
  \item \textbf{OpenAI}（2015--2017年）：作为创始成员之一，
    研究深度学习、生成模型和强化学习。
    他还曾在DeepMind（2015年暑期）实习，
    参与了深度强化学习研究。
  \item \textbf{Tesla}（2017--2022年）：担任AI高级总监，
    领导Autopilot自动驾驶的计算机视觉团队，
    负责数据标注、神经网络训练和在自研芯片上的部署
  \item \textbf{OpenAI}（2023--2024年）：回归后组建了
    专注于中间训练（midtraining）和合成数据生成的新团队
  \item \textbf{Eureka Labs}（2024年至今）：创立的AI教育平台，
    首个课程是LLM101n——``让我们构建一个讲故事的AI''
\end{itemize}

\subsubsection{Tesla：纯视觉方案的技术推手}
在Tesla的五年是Karpathy职业生涯中最具挑战性的一段时期，
也是他对工业界影响最深远的阶段。
他主导了一个在当时极具争议的技术决策：
\textbf{移除雷达和激光雷达，完全依赖摄像头}
实现自动驾驶。这个``纯视觉''（Vision-Only）方案
意味着Tesla的Autopilot/FSD系统仅使用
安装在车身周围的八个摄像头的视频输入，
通过深层卷积神经网络提取特征，
再用Transformer网络融合多视角信息，
最后通过递归神经网络（LSTM）在时间维度上整合。

Karpathy在2021年CVPR会议上的技术演讲
详细解释了这个架构的设计理念：
``人类只用两个眼睛就能开车。
如果摄像头足够好、神经网络足够大、
训练数据足够多，纯视觉方案是可行的。''
他还主导建设了Tesla内部超级计算机集群
（全球前五），以及大规模自动标注数据的离线流水线。

\subsubsection{AI时代的教育家}
Karpathy之所以在AI社区拥有独特的影响力，
不仅因为他的工程能力，更因为他是一位天才的教育者和传播者。
他的YouTube频道上关于大语言模型的深度讲解视频
（如``Deep Dive into LLMs like ChatGPT''和
``How I use LLMs''）获得了数百万次观看。

他的``Zero to Hero''编程系列从零开始构建神经网络，
被认为是互联网上最好的深度学习入门教程之一。
这个系列有一个鲜明的教学哲学：
\textbf{不使用任何库，从第一行代码开始手写一切}。
从微积分和反向传播开始，到手写一个完整的GPT模型，
整个过程完全透明。他的开源项目\texttt{nanoGPT}
和\texttt{minbpe}也遵循同样的``极简主义''原则——
用最少的代码行数实现核心功能，
让任何人都能在一台笔记本电脑上理解和运行。

\begin{keybox}[Key Quote: Karpathy的编程哲学]
``\textbf{Software 2.0}——在这个新范式中，
你不再用C++或Python编写程序。
你用神经网络的权重`编写'它们。
人类写的部分只是架构和训练代码，
而真正的`程序'是通过大规模数据优化出来的。''

——Andrej Karpathy，\textit{Software 2.0}博客文章（2017年）
\end{keybox}

2020年，他入选\textit{MIT Technology Review}
``35岁以下创新者''榜单；
2024年，他被\textit{TIME}评为AI领域
100位最具影响力人物之一。


% =============================================================
\section{GAN的发明者：Ian Goodfellow}
\label{sec:goodfellow}
% =============================================================

Ian Goodfellow于1987年出生在美国，
是深度学习领域最具创造力的年轻研究者之一。
他在斯坦福大学获得计算机科学学士和硕士学位，
最初是医学预科生，后来在Andrew Ng的课程影响下
转向了机器学习。``我意识到这会是一条
更快地影响更多领域的路径，''他回忆说。
之后他前往蒙特利尔大学，在Yoshua Bengio的指导下
攻读机器学习博士学位。

\subsection{酒吧里诞生的灵感}

2014年的一个夜晚，Goodfellow在蒙特利尔的一家酒吧
与朋友讨论技术问题时，突然想到了一个绝妙的主意：
用两个神经网络互相对抗——一个生成器（Generator）
试图创造逼真的假数据，一个判别器（Discriminator）
试图区分真假。两者在博弈中互相提升。

他当晚回家就开始编程，第一次实验就成功了。
这就是\textbf{生成对抗网络}
（Generative Adversarial Networks，GANs）的诞生。
Yann LeCun称GAN是``过去十年机器学习领域最酷的发明''。

GAN的应用范围极其广泛：从生成逼真的人脸照片
（StyleGAN）到图像超分辨率、数据增强、药物发现，
乃至后来引发争议的Deepfake视频技术。

\subsection{从OpenAI到Google到Apple再到DeepMind}
Goodfellow的职业路径经历了AI产业的核心节点：
\begin{itemize}[noitemsep]
  \item \textbf{Google}（2015--2016年，2017--2019年）：
    在Google Brain担任研究科学家和高级研究员
  \item \textbf{OpenAI}（2016--2017年）：早期研究科学家
  \item \textbf{Apple}（2019--2022年）：
    担任机器学习总监（Director of Machine Learning），
    在Apple的秘密特殊项目组工作
  \item \textbf{DeepMind}（2022年至今）：研究科学家
\end{itemize}

他还是深度学习领域最重要的教科书之一
——\textit{Deep Learning}（2016年，与Bengio和Aaron Courville合著）
的第一作者。这本书被全球数百所大学采用，
至今仍是深度学习入门的标准参考。

Goodfellow被\textit{Foreign Policy}杂志
评为100位全球思想者之一，
被\textit{MIT Technology Review}评为
35岁以下创新者之一。


% =============================================================
\section{Ilya Sutskever——从AlexNet到OpenAI到SSI}
\label{sec:sutskever}
% =============================================================

Ilya Sutskever的名字贯穿了过去十五年深度学习的几乎每一个
重大突破。他是AlexNet的合作者、序列到序列学习的发明者、
OpenAI的联合创始人和首席科学家、以及2023年``OpenAI政变''
的关键参与者。

\subsection{从俄罗斯到以色列到多伦多}

Sutskever于1986年12月8日出生在苏联高尔基市
（现俄罗斯下诺夫哥罗德）。他的家庭是犹太裔，
在苏联解体前夕的动荡中，五岁的Sutskever
随家人移居以色列耶路撒冷。
他在耶路撒冷度过了童年和少年时期，
曾就读于以色列开放大学（Open University of Israel）。
五岁时第一次看到计算机，
他就被``完全迷住了''——``I was utterly enchanted,''
他后来对《多伦多大学杂志》回忆说。
这种着迷持续到了他的青少年时代。

16岁时，Sutskever再次移民，这次来到了加拿大。
他在多伦多大学获得数学学士学位（2005年），
然后做出了改变他一生的决定：
加入Geoffrey Hinton的实验室。
在Hinton的指导下，他获得了计算机科学
硕士和博士学位，博士论文聚焦于训练递归神经网络。
Hinton后来无数次提及Sutskever的直觉天赋：
``Ilya有一种非凡的能力，能够看出什么是重要的，
而什么只是噪声。''

\subsection{关键直觉：规模就是一切}

Sutskever最重要的智识贡献可能不是任何一篇具体的论文，
而是一个\textbf{直觉}：神经网络的性能会随着
数据和模型规模的增长而持续提升。
在2011年，这还远非共识——大多数研究者认为
神经网络很快就会碰到天花板。但Sutskever说服了Krizhevsky
用更大的数据集和更深的网络来做实验，
最终催生了AlexNet。

``Ilya觉得我们应该做这件事，Alex让它跑起来了，
而我得了诺贝尔奖，''Hinton后来打趣说。

\subsection{OpenAI：从联合创始人到``政变者''}

2015年，Sutskever与Sam Altman、Greg Brockman等人
共同创立了OpenAI。作为首席科学家和联合创始人，
他领导了OpenAI的核心研究方向，
是GPT系列模型的幕后推动者。他还共同领导了
Superalignment团队——一个专注于确保超级智能AI
可以被控制和对齐的研究小组。此外，
他是AlphaGo论文的唯一Google DeepMind外部合作者，
也是CLIP、DALL-E等多个里程碑式项目的核心贡献者。
在学术影响力方面，他是历史上被引用次数最多的
计算机科学家之一，并连续三年（2022--2024年）
获得NeurIPS时间检验奖（Test of Time Award）——
一个前所未有的三连冠。

2023年11月17日，一场震惊科技界的事件发生了：
OpenAI董事会突然解雇了CEO Sam Altman。
Sutskever作为六名董事会成员之一投票支持了这一决定。
据多方报道，他的动机是出于对AI安全的深层担忧——
他认为OpenAI在商业化的道路上走得太快，
安全考量正被增长压力所边缘化。
然而，事态发展之快远超所有人的预期：
超过700名员工签署公开信威胁集体辞职，
微软（OpenAI最大投资者）施加了巨大压力。
仅仅五天后，Altman就恢复了CEO职位。

Sutskever在11月20日发文表示``深感后悔''，
称``我从未想过伤害OpenAI''。但这一事件的余波
持续了数月。Superalignment团队在他和共同领导人
Jan Leike先后离开后被解散。
2024年5月，Sutskever正式离开了他参与创立的公司。

\subsection{Safe Superintelligence Inc.}

2024年6月19日，Sutskever在X上发布了一条简洁的声明：
``我正在创办一家新公司。我们将以一条直线追求安全的超级智能，
只有一个焦点、一个目标、一个产品。''
\textbf{Safe Superintelligence Inc.}（SSI）
由他与Daniel Gross和Daniel Levy共同创立，
总部设在帕洛阿尔托和特拉维夫。

SSI的独特之处在于它的组织原则：
\textbf{在开发任何产品之前先解决安全问题}。
公司没有收入目标，没有产品路线图，
没有商业化压力——它故意将自己隔离在
硅谷的``快速发布、快速迭代''文化之外。
``SSI是我们的使命、我们的名称、
也是我们的全部产品路线图，''Sutskever写道，
``因为它是我们的唯一焦点。''

尽管没有任何产品，SSI的估值飙升速度令人瞠目。
2024年9月获得约10亿美元融资后，
2025年3月又完成了20亿美元的新一轮融资，
估值达到约300亿美元——
这使得SSI成为了历史上估值最高的
``没有产品、没有收入''的创业公司。
投资者押注的完全是Sutskever的个人智慧和判断力。

Sutskever于2022年当选为英国皇家学会院士（FRS），
2026年获得美国国家科学院工业应用科学奖。
他的故事是AI时代最深刻的道德寓言之一：
一个亲手建造了世界上最强大AI系统的人，
最终选择将毕生精力投入到确保这些系统不会伤害人类。

\begin{whybox}[Why This Matters: 数据的顿悟]
Sutskever最重要的智识遗产或许可以用一句话概括：
\textbf{``数据是AI的心脏。''}
在2011年，大多数研究者认为神经网络的性能
会很快碰到天花板。但Sutskever有一个近乎信仰的直觉：
如果你给模型更多的数据、更大的规模，
它的表现会持续提升——不会有明显的``瓶颈''。
这个直觉催生了AlexNet，推动了GPT系列的Scaling Law，
并最终塑造了整个AI产业``越大越好''的范式。
然而，讽刺的是，创造了Scaling Law信仰的人，
后来也是第一个说``仅靠扩展是不够的''的人——
这正是他创立SSI的原因。
\end{whybox}


% =============================================================
\section{DeepMind：从游戏到诺贝尔奖}
\label{sec:deepmind}
% =============================================================

% -------------------------------------------------------------
\subsection{Demis Hassabis——国际象棋神童、游戏设计师、诺贝尔奖得主}
\label{sec:hassabis}
% -------------------------------------------------------------

Demis Hassabis于1976年7月27日出生在英国伦敦北部，
父亲是希腊裔塞浦路斯人，母亲是新加坡华人。
他的人生轨迹读起来像是一部不可思议的小说。

\subsubsection{神童年代}
四岁时，Hassabis看到父亲和叔叔下国际象棋，
要求他们教他。他很快就在棋局上击败了两位长辈。
到13岁时，他已是国际象棋大师（Chess Master），
Elo等级分一度达到约2300分，
是当时英格兰同年龄段排名第二的棋手。
他在国际比赛中与成年人对弈，
展现出远超年龄的战略思维能力。
八岁时，他得到了人生中第一台计算机，
立刻用它编写了一个能下奥赛罗棋（Othello）的程序。

然而，他的野心远不止于棋盘。他很早就意识到，
比赢得棋局更有趣的是\textbf{理解游戏背后的思维机制}，
并用计算机重现它。``下棋让我开始思考：
我自己的思维过程是怎样的？
我能不能改进它？''——这个反思直接将他引向了人工智能。

\subsubsection{游戏产业与学术之路}
Hassabis的游戏生涯开始得惊人地早。
16岁时，他就加入了Peter Molyneux创立的
Bullfrog Productions，参与设计了赛博朋克策略游戏
\textit{Syndicate}。17岁时，他担任了
\textit{Theme Park}的\textbf{首席程序员}——
这是一款模拟经营游乐园的游戏，
赢得了当年的金摇杆奖（Golden Joystick Award），
全球销量数百万份，
催生了整个``经营模拟''（Tycoon）游戏品类。
考虑到他当时只有17岁，这个成就几乎不可思议。

Peter Molyneux回忆Hassabis第一天来上班的场景：
``他看起来像《指环王》里的精灵……
这个瘦小的孩子走进来，你在街上可能看都不会看一眼。
但他眼里有一种闪光：智慧的闪光。''

Hassabis随后在剑桥大学皇后学院（Queens' College）
以``双重一等''（Double First）的成绩获得计算机科学学士学位（1997年）。
毕业后他先在Peter Molyneux的新公司Lionhead Studios工作，
又创立了自己的游戏公司Elixir Studios（2000--2005年），
开发了即时战略游戏\textit{Republic: The Revolution}
和\textit{Evil Genius}。

然而，游戏产业没有完全满足他的求知欲。
2005年，他做出了一个令人惊讶的决定：
放弃商业，进入伦敦大学学院（UCL）攻读认知神经科学博士学位。

\subsubsection{神经科学：理解大脑如何想象}
Hassabis的博士研究聚焦于一个迷人的问题：
\textbf{人类大脑是如何构建对未来场景的想象的？}
他发现，海马体（hippocampus）——
一个传统上被认为只与记忆有关的脑区——
在``想象从未发生过的新场景''时同样发挥关键作用。
海马体受损的患者不仅会失去记忆，
还会失去想象新场景的能力。

这项研究发表在2007年的\textit{Proceedings of the
National Academy of Sciences}上，
随后在\textit{Science}的拓展研究中得到进一步验证，
被誉为认知神经科学领域十年来最重要的发现之一。
更关键的是，它为Hassabis后来创立DeepMind提供了
核心的生物学启示：\textbf{记忆、想象和规划
共享相同的神经计算基础}——
如果你能理解这种机制，
你就找到了通往通用智能的生物学线索。

\subsubsection{DeepMind的创立与收购}
2010年，Hassabis与Shane Legg和Mustafa Suleyman
共同创立了\textbf{DeepMind Technologies}。
公司的使命直截了当：``解决智能，然后用智能解决一切''
（Solve intelligence, then use that to solve everything else）。

2014年1月，Google以约5亿英镑收购了DeepMind
——当时AI创业公司的最高收购价。
收购后，DeepMind保持了相当大的研究自主权。

\subsubsection{AlphaGo：AI的``登月时刻''}
2016年3月，DeepMind的\textbf{AlphaGo}在首尔以4:1
击败了围棋世界冠军李世石——
这是AI历史上最具象征意义的时刻之一。
全球超过2亿人在线观看了这场比赛。
围棋因其天文数字般的可能局面（约$10^{170}$），
被认为是``AI的终极挑战''——
大多数专家预测AI至少还需要十年才能在围棋上击败人类。

这场比赛的决定性时刻是第二局中的\textbf{第37手}（Move 37）。
这是一步如此非传统的棋着，以至于职业棋手评论员
最初认为它是一个错误——``这不是人类会下的棋。''
但它被证明是决定性的一步。
Hassabis后来在2026年3月AlphaGo十周年回顾中写道：
``Move 37展示了AI的潜力，
它标志着我们现在拥有了
开始处理真实世界科学问题的技术手段。''

更令人动容的是第四局。李世石在前三局均告负后，
在第四局中下出了被称为``God's Move''的第78手，
击败了AlphaGo。当这一步棋落下时，
观赛室里的DeepMind工程师们竟然鼓起了掌——
他们为对手的创造力而喝彩。
这一刻被纪录片\textit{AlphaGo}（2017年）完美捕捉，
成为了AI历史上最具人文温度的影像记录之一。

AlphaGo的技术核心——结合深度神经网络与蒙特卡洛树搜索的
强化学习方法——主要由David Silver领导的团队开发
（见\ref{sec:david-silver}节）。

\subsubsection{AlphaFold与诺贝尔奖}
如果说AlphaGo证明了AI可以在游戏中超越人类，
那么\textbf{AlphaFold}则证明了AI可以解决
人类无法解决的科学问题。

蛋白质折叠问题——根据氨基酸序列预测蛋白质的三维结构——
是生物学领域持续了50年的大挑战。
2020年，AlphaFold 2在CASP14竞赛中以碾压性的优势获胜，
准确度达到了实验方法的水平。到2024年，
AlphaFold数据库已免费发布了超过2.14亿个蛋白质结构预测，
被来自190个国家的超过200万名研究者使用。

2024年10月9日，瑞典皇家科学院宣布将诺贝尔化学奖
授予Demis Hassabis和John Jumper（因蛋白质结构预测），
以及David Baker（因计算蛋白质设计）。
这使得Hassabis成为了从国际象棋神童到游戏设计师
到神经科学家再到诺贝尔奖得主的``不可能之路''的终点
——或者说，新的起点。
2024年他被英国王室授予爵士爵位（knighted），
2025年他入选\textit{TIME} 100最具影响力人物，
并被选为\textit{TIME}年度人物``AI建筑师''群体的代表之一。

\subsubsection{管理风格与``谨慎乐观主义''}
Hassabis对AI的态度可以被概括为\textbf{``谨慎乐观主义''}
（cautious optimism）——
这与Hinton的``悲观警报主义''和LeCun的``务实乐观主义''
形成了有趣的三角关系。
``我终生致力于这项工作，因为我相信它将成为
人类最有益的技术，''他在获得诺贝尔奖后说，
``但伴随着如此强大和变革性的技术，风险也随之而来。''

作为DeepMind的CEO，Hassabis以一种罕见的管理风格著称：
他同时是科学家和管理者，能够深入到具体的技术细节中，
又能把握组织的整体战略方向。
他还创立了\textbf{Isomorphic Labs}（2021年），
一家将AI应用于药物发现的独立公司，
利用AlphaFold的技术积累来加速新药开发。
截至2026年，他继续担任Google DeepMind的CEO
和英国政府的AI顾问。

\begin{whybox}[Why This Matters: 跨学科的力量]
Hassabis的故事是跨学科思维力量的极致展现。
国际象棋训练了他的战略思维，游戏设计教会了他如何
构建复杂系统，神经科学给了他理解智能的生物学视角。
正是这些看似不相关的经历的融合，
使他能够设想并实现像DeepMind这样的组织。
这提醒我们：\textbf{最伟大的AI突破往往
来自于AI之外的灵感。}
\end{whybox}

\bigskip

% -------------------------------------------------------------
\subsection{David Silver——AlphaGo的架构师}
\label{sec:david-silver}
% -------------------------------------------------------------

David Silver于1976年出生在英国，
在剑桥大学基督学院（Christ's College）获得学士学位（1997年），
期间与Demis Hassabis成为了朋友。
毕业后，他没有直接进入学术界，
而是共同创立了一家视频游戏公司Elixir Studios。

2004年，Silver重返学术界，
在加拿大阿尔伯塔大学攻读强化学习方向的博士学位（2009年毕业）。
他的博士论文``Reinforcement Learning and Simulation-Based
Search in Computer Go''为后来的AlphaGo奠定了技术基础。

Silver在DeepMind领导或共同领导了一系列里程碑式的项目：
\begin{itemize}[noitemsep]
  \item 用深度强化学习玩Atari游戏（\textit{Nature}，2015年）
  \item \textbf{AlphaGo}：击败围棋世界冠军（\textit{Nature}，2016年）
  \item \textbf{AlphaZero}：从零开始学习下棋、将棋和围棋，
    击败所有最强AI程序（\textit{Nature}，2017年；\textit{Science}，2018年）
  \item \textbf{MuZero}：不需要知道规则就能学会游戏
    （\textit{Science}，2020年）
  \item 对AlphaFold（蛋白质折叠）和AlphaProof
    （在国际数学奥林匹克中获得奖牌，2024年）的贡献
\end{itemize}

Silver的工作获得了ACM计算奖（2019年）、
Mensa基金会奖、英国皇家工程学院银质奖章、
以及ACM Fellow和英国皇家学会院士（FRS）等荣誉。

2025年，Silver离开DeepMind，创立了\textbf{Ineffable Intelligence}，
继续追求用强化学习构建超级智能代理的研究。

\bigskip

% -------------------------------------------------------------
\subsection{John Jumper——解开蛋白质密码的人}
\label{sec:jumper}
% -------------------------------------------------------------

John Michael Jumper于1985年出生在美国阿肯色州小石城。
他在范德堡大学（Vanderbilt University）获得物理学学士学位，
在剑桥大学获得MPhil（Marshall奖学金），
在芝加哥大学获得硕士和博士学位（2017年，
导师Tobin Sosnick和Karl Freed）。
他的博士论文题为``New Methods Using Rigorous Machine Learning
for Coarse-Grained Protein Folding and Dynamics''
——精确预见了他后来在DeepMind的工作方向。

加入DeepMind后，Jumper领导了AlphaFold 2的技术开发。
这个系统在2020年的CASP14竞赛中取得了中位GDT分数92.4
的惊人成绩（满分100，实验方法通常在90分左右），
解决了一个困扰生物学家半个世纪的问题。

2024年，38岁的Jumper与Hassabis共同获得诺贝尔化学奖
——他是历史上最年轻的化学奖得主之一。
``这太不可思议了，''他在获奖后说。
``我一直是一个计算生物学家，
我喜欢在演讲中说：我们需要这个方法行得通。
我们需要计算来解决生物学的问题——
而它真的开始行得通了。''

Jumper的其他荣誉包括\textit{Nature}年度十大科学人物（2021年）、
BBVA基金会前沿知识奖（2022年）和
Breakthrough Prize生命科学奖（2023年）。
截至2025年，他继续在Google DeepMind担任总监。


% =============================================================
\section{理论奠基者与AI哲人}
\label{sec:theorists}
% =============================================================

深度学习的爆发离不开更广泛的AI理论基础。
本节介绍几位虽然不总被归入``深度学习''阵营，
但对整个AI领域具有根本性影响的理论家。

% -------------------------------------------------------------
\subsection{Judea Pearl——因果推理的先知}
\label{sec:pearl}
% -------------------------------------------------------------

Judea Pearl于1936年9月4日出生在英属巴勒斯坦托管地的
特拉维夫（现以色列）。他在以色列理工学院（Technion）
获得电气工程学士学位（1960年），
在纽瓦克工程学院（现新泽西理工学院）获得电子学硕士学位（1961年），
在罗格斯大学获得物理学硕士学位，
在布鲁克林理工学院（现纽约大学工程学院）
获得电气工程博士学位（1965年）。

\subsubsection{贝叶斯网络与因果推理}
1969年，Pearl加入加州大学洛杉矶分校（UCLA）计算机科学系，
此后再未离开。他的两项核心贡献彻底改变了AI处理不确定性的方式：

\begin{itemize}[noitemsep]
  \item \textbf{贝叶斯网络}（1980年代）：Pearl开发了一种
    用有向无环图表示变量之间概率依赖关系的框架，
    以及在这些图上进行高效推理的信念传播算法。
    这一框架使得AI系统能够在不确定条件下进行合理推理，
    成为医学诊断、语音识别和机器人技术等领域的基石。
  \item \textbf{因果推理}（1990年代至今）：Pearl提出了
    一套基于结构因果模型的因果推理数学框架——
    do-calculus——使得从观测数据中
    推断因果关系成为可能。他的著作
    \textit{Causality: Models, Reasoning and Inference}
    （2000年）成为了统计学、经济学、流行病学等
    多个学科的标准参考。
\end{itemize}

2011年，Pearl获得图灵奖，
``以表彰他通过发展概率和因果推理的演算
对人工智能的根本贡献''。
2025年，他被选为英国皇家学会的外籍院士。

Pearl也因其个人悲剧而为公众所知：
他的儿子Daniel Pearl是《华尔街日报》记者，
2002年在巴基斯坦被基地组织关联的恐怖分子绑架并杀害。
Pearl此后一直积极参与Daniel Pearl基金会的工作。

\begin{corebox}[Core Insight: 因果推理的AI意义]
Pearl认为当前的深度学习系统存在根本性局限：
它们只能发现\textbf{相关性}，而无法理解\textbf{因果关系}。
``我们即将见证一场小型革命，''他说。
``将因果模型引入机器学习，将产生一台能理解
自己为什么这样做、哪里做错了、
如何不再重蹈覆辙、如何向用户解释自己的错误、
并能说出`对不起'且\textbf{真心如此}的机器。''
这一观点对理解当前大语言模型的局限性
和未来的改进方向至关重要。
\end{corebox}

\bigskip

% -------------------------------------------------------------
\subsection{Michael I. Jordan——机器学习的``乔丹''}
\label{sec:michael-jordan}
% -------------------------------------------------------------

Michael Irwin Jordan于1956年2月25日出生在美国
马里兰州阿伯丁。他在路易斯安那州立大学获得心理学学士学位（1978年），
在亚利桑那州立大学获得数学（统计学）硕士学位（1980年），
在加州大学圣迭戈分校获得认知科学博士学位（1985年）。

2016年，\textit{Science}杂志将Jordan评为
\textbf{全球最具影响力的计算机科学家}——
一个令许多人感到意外的排名，因为Jordan的名字
在公众中远不如Hinton或LeCun知名。
但在学术界，他的影响力是无与伦比的。

\subsubsection{学术贡献}
Jordan的贡献横跨机器学习、统计学、认知科学和
生物信息学：
\begin{itemize}[noitemsep]
  \item \textbf{潜在狄利克雷分配}（LDA，2003年，与David Blei
    和Andrew Ng合作）：一种广泛使用的主题模型
  \item 变分推理的理论框架
  \item 循环神经网络（Jordan网络）的早期形式
  \item 机器学习中概率方法的系统发展
\end{itemize}

他在MIT（1988--1998年）和UC Berkeley（1998年至今）
培养了一批深刻影响AI领域的学生，
包括Andrew Ng、Zoubin Ghahramani（剑桥教授、
Google DeepMind研究副总裁）、
David Blei（哥伦比亚大学教授）和Percy Liang等。

Jordan是美国国家科学院院士、美国国家工程院院士、
美国艺术与科学院院士、英国皇家学会外籍院士。
他获得了2022年首届世界顶尖科学家协会（WLA）奖
（计算机科学类），2020年IEEE John von Neumann Medal，
2016年IJCAI卓越研究奖，2015年David Rumelhart奖等荣誉。

Jordan曾多次发表文章指出，人们对``AI''这个术语的
过度使用和炒作可能会造成误导。他强调，
当前的系统本质上是\textbf{统计模式匹配}，
距离真正的``智能''还有很长的路要走。

\bigskip

% -------------------------------------------------------------
\subsection{Stuart Russell与Peter Norvig——AI教科书的双星}
\label{sec:russell-norvig}
% -------------------------------------------------------------

如果说AI领域有一本``圣经''，那就是
\textit{Artificial Intelligence: A Modern Approach}
（AIMA，《人工智能：一种现代方法》）。
这本由Stuart Russell和Peter Norvig合著的教科书
自1995年首版以来，已在135个国家的超过1500所大学中使用，
Google Scholar引用超过59000次，
是世界上使用最广泛的AI教科书。
第四版于2020年出版。

\subsubsection{Stuart Russell}
Stuart Jonathan Russell于1962年出生在英国朴茨茅斯。
他在牛津大学Wadham学院获得物理学一等荣誉学士学位（1982年），
在斯坦福大学获得计算机科学博士学位（1986年）。
博士毕业后即加入UC Berkeley，
目前担任杰出教授（Distinguished Professor），
持有Smith-Zadeh工程讲席。

Russell的学术贡献涵盖AI的多个核心领域，
但大约从2013年开始，他将越来越多的精力
转向了\textbf{AI安全}。他创立了UC Berkeley的
人类兼容AI中心（CHAI），并于2019年出版了
极具影响力的著作\textit{Human Compatible}
（《与人类兼容》），提出了一种基于逆强化学习
（Inverse Reinforcement Learning）的AI安全框架——
让AI系统通过观察人类行为来学习人类的目标和偏好，
而不是被赋予固定的目标函数。

Russell目前是世界经济论坛AI理事会联合主席、
OECD AI未来专家组联合主席，
并担任英国、美国和欧盟多个AI政策咨询委员会的成员。
2021年，他被BBC邀请发表里思讲座（Reith Lectures）
——英国最负盛名的公共演讲系列之一。

\subsubsection{Peter Norvig}
Peter Norvig是美国计算机科学家，
在UC Berkeley获得博士学位（1986年），
现任斯坦福大学HAI的杰出教育院士
（Distinguished Education Fellow），
同时也是Google的研究员。

Norvig的职业生涯横跨学术界、NASA和Google：
他曾在USC、斯坦福和Berkeley任教，
在NASA Ames研究中心担任计算科学部门负责人
（获得NASA杰出成就奖），
2001年加入Google后先后领导了核心搜索算法团队
和研究部门。他是AAAI Fellow、ACM Fellow、
加州科学院院士和美国艺术与科学院院士。

与Russell共同教授的一门AI在线课程
吸引了160000名学生注册，
是MOOC革命的先驱之一。


% =============================================================
\section{未来学家与AI治理先驱}
\label{sec:futurists-governance}
% =============================================================

% -------------------------------------------------------------
\subsection{Ray Kurzweil——奇点的预言者}
\label{sec:kurzweil}
% -------------------------------------------------------------

Raymond Kurzweil于1948年2月12日出生在美国纽约皇后区。
他在MIT获得计算机科学学士学位。
他是发明家、作家、企业家和未来学家——
在AI的学术界和产业界之外，
他可能是公众认知度最高的AI相关人物。

\subsubsection{发明家的履历}
Kurzweil的发明清单令人叹为观止：
\begin{itemize}[noitemsep]
  \item 第一台全字体光学字符识别（OCR）系统
  \item 第一台文本转语音阅读机（Kurzweil Reading Machine，
    第一台被盲人音乐家Stevie Wonder订购）
  \item Kurzweil K250——第一台能够逼真模拟
    大钢琴和其他管弦乐器音色的电子乐器
  \item 多项语音识别技术的突破
\end{itemize}

他因此获得了1999年美国国家技术与创新奖章
（由克林顿总统亲自颁发）、
2001年Lemelson-MIT奖（50万美元）、
21个荣誉博士学位以及三位美国总统的表彰。
PBS将他列为``塑造美国的16位革命者''之一。

\subsubsection{奇点理论}
Kurzweil的核心思想是\textbf{加速回报定律}
（The Law of Accelerating Returns）：
技术进步是指数级的，而人类的直觉是线性的，
因此我们系统性地低估了未来的变化速度。

他在2005年的著作\textit{The Singularity Is Near}
（《奇点临近》）中预测：到2029年，
AI将达到人类水平的通用智能（AGI）；
到2045年左右，人类和机器智能将合并，
达到``奇点''（Singularity）。
在2024年的续作\textit{The Singularity Is Nearer}
（《奇点更近了》）中，他重申了这些预测，
并指出AI的最新进展使得2029年的AGI时间表
正在获得越来越多的可信度。

Kurzweil自2012年起在Google工作，
目前仍然是Google的首席未来学家和AI研究员。
无论人们是否同意他的预测，
他的框架无疑深刻影响了整个AI产业对未来的想象。

\bigskip

% -------------------------------------------------------------
\subsection{Yejin Choi——常识推理的先锋}
\label{sec:yejin-choi}
% -------------------------------------------------------------

Yejin Choi是韩裔美国计算机科学家，
专注于自然语言处理、常识推理和多模态AI研究。
她目前是斯坦福大学HAI的Dieter Schwarz基金会
计算机科学教授和高级研究员，
此前在华盛顿大学Allen School担任教授（2014--2023年），
同时是Allen AI研究所（AI2）的资深研究经理。

\subsubsection{常识推理：AI最难的问题}
Choi的核心研究方向是让AI系统具备\textbf{常识推理}能力——
即人类在日常生活中不假思索就能运用的知识：
``如果你把冰淇淋放在太阳下，它会融化''、
``人们通常在生日时感到高兴''。
这些对人类来说极其简单的推理，
对AI系统来说却异常困难。

她领导开发的项目包括：
\begin{itemize}[noitemsep]
  \item \textbf{ATOMIC}：大规模常识知识图谱
  \item \textbf{COMET}：能够自动生成常识推理的神经网络模型
  \item 对大语言模型``隐含偏见''和
    ``道德推理能力''的系统性研究
\end{itemize}

\subsubsection{``天才''基金与公共影响力}
2022年10月，Choi被授予\textbf{MacArthur Fellow}
（麦克阿瑟``天才''基金），获得80万美元无条件奖金，
以表彰她``使用自然语言处理开发能够推理并理解
人类语言隐含含义的AI系统''的工作。

2023年，她入选\textit{TIME} 100 AI最具影响力人物榜单；
2025年再次入选。她是AI领域少数同时在
学术研究和公共政策讨论中都具有重要影响力的女性科学家之一。

\bigskip

% -------------------------------------------------------------
\subsection{Percy Liang——基准测试与开源的守护者}
\label{sec:percy-liang}
% -------------------------------------------------------------

Percy Liang是斯坦福大学计算机科学副教授，
斯坦福基础模型研究中心（CRFM）主任，
也是Stanford HAI的高级研究员。
他在MIT获得学士和工程硕士学位，
在UC Berkeley获得博士学位（2011年，
导师为Michael I. Jordan和Dan Klein）。

Liang的核心贡献在于为AI/ML领域建立了
\textbf{严格的评估和基准测试框架}：
\begin{itemize}[noitemsep]
  \item \textbf{HELM}（Holistic Evaluation of Language Models）：
    一个全面评估大语言模型的开源基准，
    涵盖准确性、校准度、鲁棒性、公平性和效率等多个维度。
    HELM已成为评估和比较不同LLM的行业标准之一。
  \item \textbf{CodaLab Worksheets}：促进可复现研究的平台
  \item 基础模型透明度指数和AI政策建议
\end{itemize}

他的奖项包括总统早期职业科学家和工程师奖（PECASE，2019年）、
IJCAI计算机与思想奖（2016年）、
NSF CAREER奖和Sloan研究基金等。

Liang同时也是\textbf{Together AI}的联合创始人——
一家提供开源AI基础设施的公司。
他的工作体现了一种理念：
\textbf{AI的进步不仅需要更强的模型，
更需要更好的评估方法和更高的透明度。}


% =============================================================
\section{本章小结}
\label{sec:pioneers-summary}
% =============================================================

回顾本章介绍的三十余位深度学习先驱，
几个贯穿性的主题清晰浮现：

\begin{enumerate}[noitemsep]
  \item \textbf{长期主义的回报}：Hinton、LeCun和Bengio
    在神经网络的``寒冬''中坚持了近三十年。
    Hochreiter和Schmidhuber的LSTM等了近二十年
    才获得大规模商业应用。科学上最有价值的方向
    往往需要最长的等待。

  \item \textbf{跨学科的力量}：本章中最成功的人物
    几乎都具有跨学科背景——Hinton从心理学出发，
    Jordan从认知科学出发，Hassabis融合了
    游戏设计和神经科学，Jumper从物理学和化学出发。
    深度学习本身就是数学、计算机科学、
    统计学和神经科学交叉的产物。

  \item \textbf{工程与理论的双轮驱动}：
    每一个重大突破都需要理论洞察和工程实现的结合。
    Transformer论文的八位作者既是研究者也是工程师；
    AlexNet的成功依赖于Krizhevsky将CNN
    高效实现在GPU上的工程能力。

  \item \textbf{责任意识的觉醒}：从Hinton离开Google到
    Bengio创立LawZero，从Sutskever创立SSI到
    Russell出版\textit{Human Compatible}，
    越来越多的深度学习先驱正在将精力
    从``让AI更强大''转向``让AI更安全''。
    这种转变不是偶然的——
    恰恰是最了解这项技术的人最清楚其潜在风险。

  \item \textbf{从学术到创业的大迁徙}：
    Transformer论文的八位作者全部离开Google创业；
    LeCun在65岁离开Meta创业；
    David Silver离开DeepMind创立Ineffable Intelligence。
    这反映了一个更广泛的趋势：
    AI人才正在从大公司流向创业公司，
    寻找更大的自主权和影响力。
    关于这一趋势的详细分析，
    读者可参阅姊妹篇\textit{AI Startup Landscape}。
\end{enumerate}

\begin{whybox}[Why This Chapter Matters]
理解这些先驱的故事不仅是满足历史好奇心。
在一个AI正在深刻重塑社会的时代，
了解\textbf{谁}在做出关键决策、
他们的\textbf{动机}和\textbf{价值观}是什么、
以及他们过去的\textbf{选择}如何塑造了今天的AI景观，
对于每一个关心AI未来的人来说都至关重要。
技术不是在真空中发展的——
它是由具体的人在具体的环境中创造的。
本章讲述的就是这些人的故事。
\end{whybox}
